% ============================================================
% Model: 全同粒子统计与交换对称性
% ============================================================

\section{全同粒子统计与交换对称性}
\label{sec:identical-particles}

\begin{tipbox}[关键词标签]
全同粒子 · 交换对称性 · 玻色子 · 费米子 · 泡利不相容原理 · Slater行列式 · 复合粒子统计 · 玻色-爱因斯坦凝聚
\end{tipbox}

% --------------------------------------------------------
% 1. 模型定义框
% --------------------------------------------------------
\begin{modelbox}[模型：多粒子系统的交换对称性]
\textbf{物理情境}：自然界中存在两类全同粒子——玻色子（整数自旋）和费米子（半整数自旋）。交换两个全同粒子时，玻色子体系的总波函数不变（对称），费米子体系的总波函数变号（反对称）。这一基本性质决定了多粒子系统的基态结构和统计行为。

\textbf{交换算符}：定义 $P_{12}$ 为交换粒子 1 和 2 的算符：
\[
P_{12}\Psi(\mathbf{q}_1,\mathbf{q}_2)=\Psi(\mathbf{q}_2,\mathbf{q}_1).
\]
\begin{itemize}[nosep]
    \item 玻色子：$P_{12}\Psi=+\Psi$（对称）
    \item 费米子：$P_{12}\Psi=-\Psi$（反对称）
    \item $P_{12}^2=I$，故本征值只能是 $\pm1$
\end{itemize}

\textbf{总波函数的分解}（忽略自旋-轨道耦合时）：
\[
\Psi_{\text{total}}=\Psi_{\text{space}}\otimes\chi_{\text{spin}}.
\]
\begin{itemize}[nosep]
    \item 费米子：$\Psi_{\text{total}}$ 反对称 $\Rightarrow$ 空间对称则自旋反对称，空间反对称则自旋对称
    \item 玻色子：$\Psi_{\text{total}}$ 对称 $\Rightarrow$ 空间与自旋同对称性
\end{itemize}

\textbf{泡利不相容原理}：两个全同费米子不能处于完全相同的单粒子态。若 $\psi_i=\psi_j$，则 Slater 行列式有两行相同，行列式为零。
\end{modelbox}

% --------------------------------------------------------
% 2. 关键公式框
% --------------------------------------------------------
\begin{formulabox}[核心公式]
\begin{enumerate}[label=\textbf{(\arabic*)}]
    \item \textbf{两粒子对称化/反对称化波函数}：
    \[
    \Psi_S(\mathbf{q}_1,\mathbf{q}_2)=\frac{1}{\sqrt{2}}\bigl[\psi_a(\mathbf{q}_1)\psi_b(\mathbf{q}_2)+\psi_b(\mathbf{q}_1)\psi_a(\mathbf{q}_2)\bigr],
    \]
    \[
    \Psi_A(\mathbf{q}_1,\mathbf{q}_2)=\frac{1}{\sqrt{2}}\bigl[\psi_a(\mathbf{q}_1)\psi_b(\mathbf{q}_2)-\psi_b(\mathbf{q}_1)\psi_a(\mathbf{q}_2)\bigr].
    \]
    \texteq{$S$ 对应玻色子，$A$ 对应费米子（当 $a=b$ 时 $\Psi_A=0$）}

    \item \textbf{Slater 行列式}（$N$ 费米子）：
    \[
    \Psi_A=\frac{1}{\sqrt{N!}}
    \begin{vmatrix}
    \psi_1(\mathbf{q}_1)&\psi_1(\mathbf{q}_2)&\cdots&\psi_1(\mathbf{q}_N)\\
    \psi_2(\mathbf{q}_1)&\psi_2(\mathbf{q}_2)&\cdots&\psi_2(\mathbf{q}_N)\\
    \vdots&\vdots&\ddots&\vdots\\
    \psi_N(\mathbf{q}_1)&\psi_N(\mathbf{q}_2)&\cdots&\psi_N(\mathbf{q}_N)
    \end{vmatrix}.
    \]
    \texteq{自动满足反对称性；两行相同则 $\Psi_A=0$（泡利原理）}

    \item \textbf{无限深势阱中的两粒子基态能量}（单粒子能级 $E_n=n^2\varepsilon_0$）：
    \begin{center}
    \begin{tabular}{lll}
    \toprule
    \textbf{粒子类型} & \textbf{基态占据} & \textbf{总能量} \\
    \midrule
    非全同（可区分） & $(1,1)$ & $2\varepsilon_0$ \\
    玻色子 & $(1,1)$ & $2\varepsilon_0$ \\
    费米子（自旋平行） & $(1,2)$ & $5\varepsilon_0$ \\
    费米子（自旋反平行） & $(1,1)$ & $2\varepsilon_0$ \\
    \bottomrule
    \end{tabular}
    \end{center}

    \item \textbf{复合粒子统计规则}：由 $N$ 个费米子组成的复合粒子：
    \[
    \text{整体统计}=(-1)^N\times\begin{cases}+1&\text{玻色子},\\-1&\text{费米子}.\end{cases}
    \]
    \texteq{$N$ 偶 $\Rightarrow$ 玻色子；$N$ 奇 $\Rightarrow$ 费米子}

    \item \textbf{氦同位素对比}：
    \begin{itemize}[nosep]
        \item $^4$He：2p+2n=4 费米子（偶）$\Rightarrow$ \textbf{玻色子}，极低温可发生玻色-爱因斯坦凝聚和超流
        \item $^3$He：2p+1n=3 费米子（奇）$\Rightarrow$ \textbf{费米子}，极低温形成费米液体，需 Cooper 配对才能超流
    \end{itemize}
\end{enumerate}
\end{formulabox}

% --------------------------------------------------------
% 3. 训练点框
% --------------------------------------------------------
\begin{drillbox}[训练点与考法变体]
\trainingpoint{1} \textbf{波函数对称性判定}：给定两粒子波函数 $\Psi(x_1,x_2)$，判断交换 $x_1\leftrightarrow x_2$ 后的行为，确定其描述的是玻色子还是费米子。

\trainingpoint{2} \textbf{势阱中多粒子基态构造}：给定势阱（无限深/谐振子/库仑势）和粒子数 $N$，分别写出玻色子和费米子的基态波函数（空间部分）和总能量。

\trainingpoint{3} \textbf{费米子基态与费米能}：$N$ 个自旋-$1/2$ 费米子在三维无限深势阱中，基态填满到费米能级 $E_F$。计算费米波矢 $k_F=(3\pi^2n)^{1/3}$ 和费米能 $E_F=\hbar^2k_F^2/(2m)$。

\trainingpoint{4} \textbf{复合粒子统计判定}：给定原子核组成（质子数+中子数），判定该原子是玻色子还是费米子。例如：$^{12}$C（6p+6n=12，玻色子）、$^{13}$C（6p+7n=13，费米子）。

\trainingpoint{5} \textbf{自旋-空间波函数耦合}：两个自旋-$1/2$ 费米子，总自旋 $S=0$（单态，反对称）或 $S=1$（三重态，对称）。若空间波函数对称（如 $\phi_1(x_1)\phi_1(x_2)$），则必须配单态；若空间反对称，则配三重态。

\trainingpoint{6} \textbf{交换能（Fock 项）}：由于波函数对称化要求，全同粒子间出现额外的交换相互作用能。对于费米子（自旋平行），交换能使同自旋电子相互排斥（即使在无库仑相互作用时也有``交换孔''）。
\end{drillbox}

% --------------------------------------------------------
% 4. 解题技巧框
% --------------------------------------------------------
\begin{tipbox}[快速识别与解题技巧]
\begin{itemize}
    \item \examnote{识别标志}：题干出现``全同粒子''''玻色子''''费米子''''泡利原理''''交换对称性''''Slater行列式''，立即定位本模型。
    \item \examnote{两粒子波函数构造速查}：
    \begin{center}
    \fbox{\parbox{0.9\linewidth}{\centering
    \textbf{单粒子态 $a$ 和 $b$（$a\neq b$）}\\[4pt]
    玻色子：$\Psi_S=\frac{1}{\sqrt{2}}(|a,b\rangle+|b,a\rangle)$\\[2pt]
    费米子：$\Psi_A=\frac{1}{\sqrt{2}}(|a,b\rangle-|b,a\rangle)$\\[2pt]
    若 $a=b$：玻色子 $\Psi_S=|a,a\rangle$；费米子 $\Psi_A=0$（禁止）
    }}
    \end{center}
    \item \examnote{基态能量速算}：
    \begin{itemize}[nosep]
        \item 非全同/玻色子：全部粒子堆积到最低能级
        \item 费米子（无自旋简并）：$E=\sum_{n=1}^Nn^2\varepsilon_0=\frac{N(N+1)(2N+1)}{6}\varepsilon_0$
        \item 费米子（自旋-$1/2$，自旋简并=2）：每能级可填2个，$E=2\sum_{n=1}^{N/2}n^2\varepsilon_0$
    \end{itemize}
    \item \examnote{复合粒子统计判定口诀}：``偶数费米子成玻色，奇数费米子成费米''。质子、中子、电子都是费米子，数总数即可。
    \item \examnote{Slater 行列式速写}：$N$ 个费米子占据单粒子态 $\psi_1,\ldots,\psi_N$，行列式的第 $i$ 行是 $\psi_i$，第 $j$ 列是粒子 $j$ 的坐标。交换两粒子 $\Leftrightarrow$ 交换两列 $\Leftrightarrow$ 行列式变号。
\end{itemize}
\end{tipbox}

% --------------------------------------------------------
% 5. 易错点框
% --------------------------------------------------------
\begin{errorbox}{常见失误与陷阱}
\begin{itemize}
    \item \textbf{费米子基态能量计算错误}：两个自旋平行的费米子不能占据同一能级，基态是一个在 $n=1$、一个在 $n=2$。常见错误是忽略泡利原理，直接算 $2E_1$。
    \item \textbf{自旋简并的遗漏}：自旋-$1/2$ 费米子每个空间能级可容纳 2 个（自旋向上/向下）。若题目未说明``自旋平行''，默认可考虑自旋反平行，此时两个粒子可同占 $n=1$。
    \item \textbf{波函数归一化}：对称化/反对称化波函数 $|\Psi_{S/A}\rangle=\frac{1}{\sqrt{2}}(|a,b\rangle\pm|b,a\rangle)$ 中，$1/\sqrt{2}$ 的归一化因子不能遗漏。若 $a=b$，玻色子的归一化变为 $1$（不是 $1/\sqrt{2}$）。
    \item \textbf{交换对称性 vs 空间对称性}：费米子的``总波函数反对称''不等于``空间波函数反对称''。如果自旋部分是对称的（如三重态），空间部分必须反对称；反之亦然。
    \item \textbf{复合粒子的内部结构}：$^4$He 原子整体是玻色子，但这不意味着其内部的电子不受泡利原理约束。电子仍是费米子，必须满足泡利原理；只是整个原子（含核）在交换时表现为玻色子。
    \item \textbf{玻色-爱因斯坦凝聚的条件}：理想玻色气体的 BEC 临界温度 $T_c=\frac{2\pi\hbar^2}{mk_B}\left(\frac{n}{\zeta(3/2)}\right)^{2/3}$。需要低温 + 高密度的组合，不是单纯低温即可。
\end{itemize}
\end{errorbox}

% --------------------------------------------------------
% 6. 物理图像与关联模型
% --------------------------------------------------------
\begin{insightbox}[物理图像与模型关联]
\textbf{物理图像}：

想象一个礼堂（势阱）和观众（粒子）：
\begin{itemize}[nosep]
    \item \textbf{非全同粒子}：每个观众都有独特的名字和座位偏好，大家各自找最舒服的座位（最低能级），可以挤在同一个座位上。
    \item \textbf{玻色子}：所有观众都长得一模一样，喜欢挤在一起。基态时所有人都坐在同一个最舒服的座位上——这就是玻色-爱因斯坦凝聚。
    \item \textbf{费米子}：所有观众也长得一样，但有个奇怪的规定：任何两个观众不能坐在完全相同的座位上（泡利原理）。所以他们必须分散到不同的座位上，从低到高依次就座。
    \item \textbf{自旋的作用}：如果费米子观众可以``面朝前''或``面朝后''，那么同一个座位上可以坐两个面朝不同方向的观众——这就是自旋简并。
\end{itemize}

\textbf{氦同位素的极低温行为对比}：
\begin{center}
\begin{tabular}{lll}
\toprule
\textbf{性质} & \textbf{$^4$He（玻色子）} & \textbf{$^3$He（费米子）} \\
\midrule
核组成 & 2p + 2n（4 费米子） & 2p + 1n（3 费米子） \\
统计 & 玻色子 & 费米子 \\
极低温 & 超流（BEC，$T_\lambda\approx2.17$K） & 费米液体，需配对才能超流 \\
配对机制 & 无（天然凝聚） & p 波 Cooper 配对（自旋三重态） \\
声速 & 有普通声速 + 超流声速 & 两种声速（零声 + 第一声） \\
\bottomrule
\end{tabular}
\end{center}

\textbf{与相似模型的对比}：
\begin{center}
\begin{tabular}{llll}
\toprule
\textbf{对比维度} & \textbf{Maxwell-Boltzmann} & \textbf{Bose-Einstein} & \textbf{Fermi-Dirac} \\
\midrule
分布 & $e^{-\beta(E-\mu)}$ & $\frac{1}{e^{\beta(E-\mu)}-1}$ & $\frac{1}{e^{\beta(E-\mu)}+1}$ \\
高能极限 & 相同（经典） & 相同 & 相同 \\
低温行为 & 冻结 & 凝聚到基态 & 形成费米面 \\
自旋 & 任意 & 整数 & 半整数 \\
\bottomrule
\end{tabular}
\end{center}

\textbf{推广方向}：
\begin{itemize}[nosep]
    \item 分数统计（任意子）：二维系统中自旋可为任意值，交换相位 $e^{i\theta}$，$\theta=0$（玻色子）、$\theta=\pi$（费米子）或任意值
    \item 量子霍尔效应中的复合费米子：电子绑定偶数个磁通量子，变换为复合费米子
    \item 超导体中的 Cooper 对：两个电子（费米子）配对形成玻色子，发生 BEC 导致超导
\end{itemize}
\end{insightbox}

% --------------------------------------------------------
% 7. 推导附录
% --------------------------------------------------------
\begin{derivationbox}[推导细节（自我检验用）]
\textbf{Step 1：两费米子（自旋平行）的基态}

单粒子能级：$E_n=n^2\varepsilon_0$，$\varepsilon_0=\pi^2\hbar^2/(2ma^2)$。

自旋平行 $\Rightarrow$ 自旋部分对称 $\Rightarrow$ 空间部分必须反对称。

两个粒子不能同占 $n=1$（否则空间波函数 $\phi_1(x_1)\phi_1(x_2)$ 对称，总波函数无法反对称）。

基态：一个在 $n=1$，一个在 $n=2$。反对称空间波函数：
\[
\Psi_A(x_1,x_2)=\frac{1}{\sqrt{2}}\bigl[\phi_1(x_1)\phi_2(x_2)-\phi_2(x_1)\phi_1(x_2)\bigr].
\]

总能量：$E=E_1+E_2=\varepsilon_0+4\varepsilon_0=5\varepsilon_0$。

\textbf{Step 2：第一激发态}

下一个可用组合：$(1,3)$ 或 $(2,2)$。$(2,2)$ 空间对称，不允许（自旋平行时空间必须反对称）。故第一激发态为 $(1,3)$：
\[
E=\varepsilon_0+9\varepsilon_0=10\varepsilon_0.
\]

\textbf{Step 3：玻色子基态}

玻色子 $\Rightarrow$ 空间波函数对称。两个粒子可同占 $n=1$：
\[
\Psi_S(x_1,x_2)=\phi_1(x_1)\phi_1(x_2).
\]

总能量：$E=2E_1=2\varepsilon_0$。

\textbf{Step 4：复合粒子统计证明}

设复合粒子 A 由 $N$ 个费米子组成，复合粒子 B 也由 $N$ 个全同费米子组成。

交换 A 和 B $\Leftrightarrow$ 交换它们的所有组分。

交换两个费米子产生因子 $(-1)$。交换 $N$ 个费米子从 A 到 B 等价于 $N$ 次两两交换，总因子 $(-1)^N$。

若 $N$ 为偶数，$(-1)^N=+1$，复合粒子表现为玻色子；
若 $N$ 为奇数，$(-1)^N=-1$，复合粒子表现为费米子。

\textbf{Step 5：Slater 行列式的反对称性}

$N$ 费米子波函数：
\[
\Psi_A(\mathbf{q}_1,\ldots,\mathbf{q}_N)=\frac{1}{\sqrt{N!}}\det\bigl[\psi_i(\mathbf{q}_j)\bigr].
\]

交换粒子 $k$ 和 $l$：行列式的第 $k$ 列和第 $l$ 列互换，行列式变号。故 $\Psi_A$ 反对称。

若 $\psi_i=\psi_j$（两行相同），则 $\det=0$，$\Psi_A=0$。这就是泡利不相容原理的数学表达。
\end{derivationbox}

\vspace{1em}
